% Part: normal-modal-logic
% Chapter: filtrations
% Section: examples-of-filtrations

\documentclass[../../../include/open-logic-section]{subfiles}

\begin{document}

\olfileid{nml}{fil}{exf}

\olsection{వడపోతల ఉదాహరణలు}

వడపోతలు ఉన్నాయని ఇంకా చూపలేదు. అయితే ఏ
నమూనా~$\mModel{M}$కైనా $\Gamma$ ద్వారా
$\mModel{M}$కు చేసే అనేక
వడపోతలు ఉన్నాయి. వాటిలో ప్రత్యేకంగా రెండింటిని,
అత్యంత సూక్ష్మ వడపోతను, అత్యంత స్థూల వడపోతను
గుర్తిద్దాం. ఒకే నమూనాకు చేసే వడపోతలు వాటి ప్రాప్యత
సంబంధంలో వేరుగా ఉంటాయి (
\olref[fil]{defn:filtration}లో $W^*$,~$V^*$ ఏమిటో
నేరుగా నిర్దేశించాం). అత్యంత సూక్ష్మ వడపోతలో
సంబంధిత లోకాల జతలు సాధ్యమైనంత తక్కువగా, అత్యంత
స్థూల వడపోతలో సాధ్యమైనంత ఎక్కువగా ఉంటాయి.

\begin{defn}
  $\Gamma$ ఉపసూత్రాల పరంగా సంవృతమై ఉంటే,
  నమూనా $\mModel{M}$కు చేసే \emph{అత్యంత సూక్ష్మ}
  వడపోత $\mModel{M^*}$ను ఇలా నిర్వచిస్తాం:
  \[
  R^*[u][v] \quad \text{అప్పుడు మరియు అప్పుడు మాత్రమే} \quad \exists u'\in [u] \;
  \exists v' \in [v] : Ru'v'.
  \]
\end{defn}

\begin{prop}\ollabel{prop:finest}
  అత్యంత సూక్ష్మ వడపోత $\mModel{M^*}$ నిజంగానే వడపోత.
\end{prop}

\begin{proof}
  ఇలా నిర్వచించిన $R^*$,
  \olref[fil]{defn:filtration}\olref[fil]{defn:filtration-R}ను
  నెరవేరుస్తుందో తనిఖీ చేయాలి. మూడు షరతులనూ వరుసగా
  చూద్దాం.

  $Ruv$ అయితే $u \in [u]$, $v \in [v]$ కనుక
  $R^*[u][v]$ కూడా. అందువల్ల
  \olref[fil]{defn:filtration-R1} నెరవేరింది.

  \iftag{prvBox}{\iftag{probBox}{
      \olref[fil]{defn:filtration-R2} నిర్ధారణను
      వ్యాయామంగా వదిలేస్తాం.}{
        \olref[fil]{defn:filtration-R2} కోసం
        $\Box!A \in \Gamma$, $R^*[u][v]$,
        $\mSat{M}{\Box!A}[u]$ అనుకుందాం. $R^*$ నిర్వచనం
        ప్రకారం, $u' \equiv u$, $v' \equiv v$ అయ్యే,
        $Ru'v'$ను నెరవేర్చే ప్రతినిధులు ఉన్నాయి.
        $u$, $u'$లు $\Gamma$పై ఏకీభవిస్తాయి కనుక
        $\mSat{M}{\Box!A}[u']$ కూడా; అందువల్ల
        $\mSat{M}{!A}[v']$. $\Gamma$ ఉపసూత్రాల పరంగా
        సంవృతమై ఉంది కనుక, $v$, $v'$లు
        \tetoken{ఉపసూత్రాల}{!!{formula}s}లోని $!A$పై
        ఏకీభవిస్తాయి. కాబట్టి కావలసినట్లు
        $\mSat{M}{!A}[v]$.}}{}

  \iftag{prvDiamond}{\iftag{probDiamond}{
      \olref[fil]{defn:filtration-R3} నిర్ధారణను
      వ్యాయామంగా వదిలేస్తాం.}{
        \olref[fil]{defn:filtration-R3}ను తనిఖీ చేయడానికి
        $\Diamond!A \in \Gamma$, $R^*[u][v]$,
        $\mSat{M}{!A}[v]$ అనుకుందాం. $R^*$ నిర్వచనం
        ప్రకారం, $u' \equiv u$, $v' \equiv v$ అయ్యే,
        $Ru'v'$ను నెరవేర్చే ప్రతినిధులు ఉన్నాయి.
        $v$, $v'$లు~$\Gamma$పై ఏకీభవిస్తాయి;
        $\Gamma$ ఉపసూత్రాల పరంగా సంవృతం కనుక
        \tetoken{ఉపసూత్రాల}{!!{formula}s} సంవృతత వల్ల
        $\mSat{M}{!A}[v']$ కూడా. అందువల్ల
        $\mSat{M}{\Diamond !A}[u']$. $u$, $u'$లు కూడా
        $\Gamma$పై ఏకీభవిస్తాయి కనుక
        $\mSat{M}{\Diamond !A}[u]$.}}{}
\end{proof}

\begin{probtag}{probBox,probDiamond}
  \olref[nml][fil][exf]{prop:finest} నిరూపణను పూర్తి చేయండి.
\end{probtag}

\begin{defn}
  $\Gamma$ ఉపసూత్రాల పరంగా సంవృతమై ఉంటే,
  నమూనా~$\mModel{M}$కు చేసే \emph{అత్యంత స్థూల}
  వడపోత~$\mModel{M^*}$ను, $R^*[u][v]$ అయితే,
  అప్పుడు మాత్రమే ఈ
  \iftag{notprvBox,notprvDiamond}{కింది షరతు
    నెరవేరేలా:}{కింది \emph{రెండు} షరతులూ నెరవేరేలా:}
  నిర్వచిస్తాం.
  \begin{tagenumerate}{prvBox,prvDiamond}
  \tagitem{prvBox}{\ollabel{defn:coarsest-Box}$\Box!A \in \Gamma$,
    $\mSat{M}{\Box!A}[u]$ అయితే
    $\mSat{M}{!A}[v]$\iftag{prvDiamond}{;}{.}}{}
  \tagitem{prvDiamond}{\ollabel{defn:coarsest-Diamond}$\Diamond!A
    \in \Gamma$, $\mSat{M}{!A}[v]$ అయితే
    $\mSat{M}{\Diamond!A}[u]$.}{}
  \end{tagenumerate}
\end{defn}

\begin{prop}
  అత్యంత స్థూల వడపోత $\mModel{M^*}$ నిజంగానే వడపోత.
\end{prop}

\begin{proof}
  $R^*$ నిర్వచనం ప్రకారం, $Ruv$ నుంచి $R^*[u][v]$
  వస్తుందని మాత్రమే ఇక నిర్ధారించాలి. $Ruv$
  అనుకుందాం. \iftag{prvBox}{$\Box!A \in \Gamma$,
    $\mSat{M}{\Box!A}[u]$ అయితే, స్పష్టంగా
    $\mSat{M}{!A}[v]$\iftag{prvDiamond}{;
      \olref{defn:coarsest-Box} నెరవేరింది. }{.}}{}\iftag{prvDiamond}{$\Diamond!A \in
    \Gamma$, $\mSat{M}{!A}[v]$ అనుకుందాం. $Ruv$
    కనుక $\mSat{M}{\Diamond!A}[u]$\iftag{prvBox}{;
      \olref{defn:coarsest-Diamond} నెరవేరింది}{}.}{}
\end{proof}

\begin{ex}
  $W = \PosInt$, $Rnm$ అయితే, అప్పుడు మాత్రమే
  $m = n + 1$; $V(p) = W \cap \Setabs{2n}{n \in
    \Nat}$ అని ఉంచుదాం. నమూనా
  $\mModel{M} = \tuple{W, R, V}$ను
  \olref{fig:ex-filtration}లో చిత్రించాం. లోకాలు
  $1$, $2$ మొదలైనవి. ప్రతి లోకం తన ఉత్తరవర్తి అయిన
  ఒకే లోకానికి ప్రాప్యత కలిగి ఉంటుంది; $p$ సరి
  సంఖ్యల వద్ద మాత్రమే సత్యం.
  \sourcecorrection{OLTENMLFILEXF-002}{ఈ ఎడిషన్‌లో
    $\Nat$లో $0$ ఉంది; కానీ లోకాల సమితి $\PosInt$లో
    అది లేదు. స్థిర మూలంలోని విలువ సమితిలో $2\cdot0$
    కూడా చేరుతుంది. కాబట్టి విలువ నిర్ణయాన్ని $W$తో
    ఛేదించి లోకాలకే పరిమితం చేశాం; చిత్రంలోని
    సత్యమూల్యాలు మారవు.}

  \begin{figure}
    \centering
    \begin{tikzpicture}[modal]
      \node[world] (1) [label=below:\mFalse{p}] {$1$};
      \node[world] (2) [label=below:\mTrue{p}, right=of 1] {$2$};
      \node[world] (3) [label=below:\mFalse{p}, right=of 2] {$3$};
      \node[world] (4) [label=below:\mTrue{p}, right=of 3] {$4$};
      \node[phantom] (5) [right=of 4] {};
      \draw[->] (1) to (2);
      \draw[->] (2) to (3);
      \draw[->] (3) to (4);
      \draw[dotted] (4) to (5);
    \end{tikzpicture}

    \begin{tikzpicture}[modal]
      \node[world] (1) [label=below:\mFalse{p}] {$[1]$};
      \node[world] (2) [label=below:\mTrue{p}, right=of 1] {$[2]$};
      \draw[->,bend left] (1) to (2);
      \draw[->,bend left] (2) to (1);

      \node[world] (3) [label=below:\mFalse{p},right=of 2] {$[1]$};
      \node[world] (4) [label=below:\mTrue{p}, right=of 3] {$[2]$};
      \draw[->,bend left] (3) to (4);
      \draw[->,bend left] (4) to (3);
      \draw[->,reflexive above] (4) to (4);
    \end{tikzpicture}
    \caption{అనంత నమూనా మరియు దాని వడపోతలు.}
    \ollabel{fig:ex-filtration}
  \end{figure}

  ఇప్పుడు $\Gamma$ అనేది~$\Box p \lif p$ యొక్క
  \tetoken{ఉపసూత్రాల}{!!{formula}s} సమితి అనుకుందాం;
  అంటే $\{p, \Box p, \Box p \lif p\}$.
  $p$ సరి సంఖ్యల వద్ద మాత్రమే సత్యం; $\Box p$
  బేసి సంఖ్యల వద్ద మాత్రమే సత్యం. అందువల్ల
  $\Box p \lif p$ సరి సంఖ్యల వద్ద మాత్రమే సత్యం.
  మరోలా చెప్పాలంటే, ప్రతి బేసి సంఖ్య $\Box p$ను
  సత్యం చేస్తుంది; కానీ $p$, $\Box p \lif p$లను
  అసత్యం చేస్తుంది. ప్రతి సరి సంఖ్య $p$,
  $\Box p \lif p$లను సత్యం చేస్తుంది; కానీ
  $\Box p$ను అసత్యం చేస్తుంది. కనుక
  $W^* = \{ [1], [2] \}$; ఇక్కడ
  $[1] = \{1, 3, 5, \dots\}$,
  $[2] = \{2, 4, 6, \dots\}$. $2 \in V(p)$ కనుక
  $[2] \in V^*(p)$. మొదటి వర్గంలోని లోకాలన్నీ బేసి,
  అందువల్ల $1 \notin V(p)$ మాత్రమే కాక మొత్తం
  వర్గమే విలువ సమితికి వెలుపల ఉంది; కనుక
  $[1] \notin V^*(p)$. అందువల్ల
  $V^*(p) = \{[2]\}$.

  $W^*$ ఆధారంగా చేసే ఏ వడపోత ప్రాప్యత సంబంధంలోనైనా
  $\tuple{[1], [2]}, \tuple{[2],[1]}$ తప్పనిసరిగా
  ఉండాలి. ఎందుకంటే $R12$ కనుక
  \olref[fil]{defn:filtration}\olref[fil]{defn:filtration-R1}
  వల్ల $R^*[1][2]$ కావాలి; అలాగే $R23$ కనుక
  $R^*[2][3]$ కావాలి; $[3]=[1]$. కానీ
  $\tuple{[1],[1]}$ ఉండకూడదు: అది ఉంటే
  $R^*[1][1]$, $\mSat{M}{\Box p}[1]$ ఉండి,
  $\mSat/{M}{p}[1]$ ఉంటుంది; ఇది
  \olref[fil]{defn:filtration-R2}కు విరుద్ధం.
  $R^*[2][2]$ను ఏ షరతూ తప్పనిసరి చేయదు,
  నిషేధించదు. కాబట్టి~$\mModel{M}$కు రెండు వడపోతలు
  సాధ్యం; వాటి ప్రాప్యత సంబంధాలు:
  \[
  \{\tuple{[1],[2]}, \tuple{[2],[1]}\} \text{ మరియు }
  \{\tuple{[1],[2]}, \tuple{[2],[1]}, \tuple{[2],[2]}\}.
  \]
  రెండు సందర్భాల్లోనూ~$[1]$ వద్ద $p$,
  $\Box p \lif p$ అసత్యం, $\Box p$ సత్యం;
  $[2]$ వద్ద $p$, $\Box p \lif p$ సత్యం,
  $\Box p$ అసత్యం.
\end{ex}


\begin{prob}
  కింది నమూనా $\mModel{M} = \tuple{W, R, V}$ను
  పరిశీలించండి. ఇక్కడ $W
  = \Setabs{0\sigma}{\sigma \in \Bin^*}$,
  అంటే $0$తో ప్రారంభమయ్యే $0$,~$1$ల క్రమాల సమితి;
  $R\sigma\sigma'$ అయితే, అప్పుడు మాత్రమే
  $\sigma' = \sigma 0$ లేదా $\sigma' = \sigma 1$.
  అలాగే $V(p) = W \cap \Setabs{\sigma
    0}{\sigma \in \Bin^*}$,
  $V(q) = W \cap \Setabs{\sigma 1}{\sigma \in
    \Bin^* \setminus \{1\}}$. చిత్రం ఇలా ఉంది:
  \sourcecorrection{OLTENMLFILEXF-001}{స్థిర మూలం
    $V(p)$, $V(q)$ సమితులను $W$కు పరిమితం చేయలేదు;
    అందువల్ల కొన్ని $1$తో మొదలయ్యే క్రమాలు
    లోకాలు కాకపోయినా విలువ సమితుల్లోకి రావచ్చు.
    నమూనా విలువ నిర్ణయం $W$ ఉపసమితులనే ఇవ్వాలి
    కనుక రెండింటికీ $W$తో ఛేదనం చేర్చాం. మూలంలోని
    $\{1\}$ మినహాయింపును యథాతథంగా ఉంచాం.}
  \begin{center}
    \begin{tikzpicture}[modal,
        node distance=2cm
      ]

      \node[world] (0) [label={below:\mTrue{p}\\\mFalse{q}},font=\tiny] {$0$};
      \node[world] (00) [label={below:\mTrue{p}\\\mFalse{q}},
        above right=of 0,font=\tiny] {$00$};
      \node[world] (000) [label={below:\mTrue{p}\\\mFalse{q}},
        right=of 00, yshift=.8cm, font=\tiny] {$000$};
      \node[phantom] (0000) [right=of 000, yshift=.5cm] {};
      \node[phantom] (0001) [right=of 000, yshift=-.5cm] {};
      \node[world] (001) [label={below:\mFalse{p}\\\mTrue{q}},
        right=of 00, yshift=-.8cm, font=\tiny] {$001$};
      \node[phantom](0010) [right=of 001, yshift=.5cm] {};
      \node[phantom](0011) [right=of 001, yshift=-.5cm] {};
      \node[world] (01) [label={below:\mFalse{p}\\\mTrue{q}},
        below right=of 0,font=\tiny] {$01$};
      \node[world] (010) [label={below:\mTrue{p}\\\mFalse{q}},
        right=of 01, yshift=.8cm,font=\tiny] {$010$};
      \node[phantom] (0100) [right=of 010, yshift=.5cm] {};
      \node[phantom] (0101) [right=of 010, yshift=-.5cm] {};
      \node[world] (011) [label={below:\mFalse{p}\\\mTrue{q}},
        right=of 01, yshift=-.8cm,font=\tiny] {$011$};
      \node[phantom] (0110) [right=of 011, yshift=.5cm] {};
      \node[phantom] (0111) [right=of 011, yshift=-.5cm] {};

      \draw[->] (0) to (00);
      \draw[->] (0) to (01);
      \draw[->] (00) to (000);
      \draw[dotted] (000) to (0000);
      \draw[dotted] (000) to (0001);
      \draw[->] (00) to (001);
      \draw[dotted] (001) to (0010);
      \draw[dotted] (001) to (0011);
      \draw[->] (01) to (010);
      \draw[dotted] (010) to (0100);
      \draw[dotted] (010) to (0101);
      \draw[->] (01) to (011);
      \draw[dotted] (011) to (0110);
      \draw[dotted] (011) to (0111);
    \end{tikzpicture}
  \end{center}
  ప్రతి~$w$ వద్ద
  $\mSat/{M}{\Box(p \lor q) \lif (\Box p \lor \Box q)}[w]$.

  ఇప్పుడు $\Gamma$ అనేది~$\Box(p \lor q) \lif
  (\Box p \lor \Box q)$ యొక్క
  \tetoken{ఉపసూత్రాల}{!!{formula}s} సమితి అనుకుందాం.
  $W^*$,~$V^*$ ఏమిటి? $\mModel{M}$కు చేసే
  అత్యంత సూక్ష్మ వడపోత ప్రాప్యత సంబంధం ఏమిటి?
  అత్యంత స్థూల వడపోత ప్రాప్యత సంబంధం ఏమిటి?
\end{prob}

\end{document}
